% !TeX program = xelatex
% UTF-8; compile from this directory with XeLaTeX.
% Source: Google Drive / Notebook / congenital acyanotic diseases.
% Original photographs, in handwritten page order:
% pic01.jpg = 20261016_053211003.jpg
% pic02.jpg = 20261016_053216154.jpg
% pic03.jpg = 20261016_053219991.jpg
% pic04.jpg = 20261016_053224983.jpg
% pic05.jpg = 20261016_053228053.jpg
% pic06.jpg = 20261016_053232835.jpg
% pic07.jpg = 20261016_053235880.jpg
% pic08.jpg = 20261016_053243386.jpg
% pic09.jpg = 20261016_053246263.jpg
% pic10.jpg = 20261016_053251503.jpg
% pic11.jpg = 20261016_053254457.jpg
% Dual-mode guard: compiles standalone with XeLaTeX, or is \input from the
% XePersian notebook (main.tex) where its local preamble below is skipped.
\newif\ifCongenitalAcyanoticStandalone
\ifdefined\settextfont
  \CongenitalAcyanoticStandalonefalse
\else
  \CongenitalAcyanoticStandalonetrue
\fi
\ifCongenitalAcyanoticStandalone
\documentclass[12pt,a4paper]{article}
\usepackage[margin=22mm]{geometry}
\usepackage{amsmath,amssymb}
\usepackage{graphicx}
\usepackage{xepersian}
\IfFontExistsTF{Amiri}{\settextfont{Amiri}}{\settextfont{DejaVu Sans}}
\setlatintextfont{DejaVu Serif}
\setlength{\parindent}{0pt}
\setlength{\parskip}{5pt}
\setlength{\emergencystretch}{3em}
\newcommand{\E}[1]{\lr{#1}}
\newcommand{\M}[1]{\lr{\(#1\)}}
\newcommand{\U}[1]{\textbf{[خوانش نامطمئن: #1]}}
\newcommand{\Cut}{\textbf{[ادامه در لبهٔ عکس بریده شده است]}}
\newcommand{\Page}[1]{\clearpage\section*{تصویر #1 ـ فصل ۴۷۵}}
\begin{document}
\fi

% Helper macros used inside the body must exist in notebook mode too.
\providecommand{\E}[1]{\lr{#1}}
\providecommand{\M}[1]{\lr{\(#1\)}}
\providecommand{\U}[1]{\textbf{[خوانش نامطمئن: #1]}}
\providecommand{\Cut}{\textbf{[ادامه در لبهٔ عکس بریده شده است]}}
\providecommand{\Page}[1]{\clearpage\section*{تصویر #1 ـ فصل ۴۷۵}}

\begin{center}
{\Large\bfseries بیماری‌های مادرزادی آسیانوتیک: شانت‌های چپ به راست}\par
\E{Nelson 22}\par
فصل ۴۷۵
\end{center}
رونویسی به ترتیب ۱۱ تصویر است. فلش‌ها و شاخه‌های دست‌نویس به صورت رابطه یا فهرست آمده‌اند. عبارت‌های نامطمئن و نوشته‌های بریده‌شده در لبهٔ عکس مشخص شده‌اند. تصاویر اصلی برای مقابله در پایان همین سند آمده‌اند. سربرگ چاپی برگه‌ها «موضوع» و «تاریخ» است؛ موضوع در تصویر اول عنوان بالا و در تصاویر بعدی «۴۷۵» است. تاریخ‌ها خالی و با دو ممیز هستند. شماره‌های دست‌نویس پایین برگه‌ها به ترتیب ۱ تا ۱۱ است.

\section*{تصویر ۱}
\subsection*{۱. \E{Atrial Septal Defect (ASD)}}
\begin{itemize}
\item در هر بخشی از سپتوم می‌تواند رخ دهد، یا حتی چند نقص:
\E{Secundum}، \E{Primum}، \E{Sinus Venosus}.
\item حتی ممکن است کلاً سپتوم وجود نداشته باشد؛ \U{مولتیپل نباشد}؛ \E{Functional Single Atrium}.
\item نوع \E{Secundum} شایع‌ترین است: ۷\% تمام نقص‌های مادرزادی قلب.
\item وراثت: اغلب اسپورادیک؛ \E{AD} بخشی از سندرم \E{Holt--Oram}.
\end{itemize}
\begin{latin}
ASD + 1st degree AV block, Phocomelia, Absent thumbs
\end{latin}
\E{Phocomelia}: فقدان پروگزیمال اندام‌ها.

\subsection*{\E{PFO}}
\E{patent foramen} (عبارت حاشیه‌ای در اصل).
\begin{itemize}
\item یک \E{PFO} ایزوله یک یافتهٔ شایع در دوران شیرخوارگی است.
\item اهمیت همودینامیک ندارد.
\item یک \E{ASD} حساب نمی‌شود.
\item اگر یک آنومالی دیگر وجود داشته باشد که فشار دهلیز راست را بالا ببرد، وجود \E{PFO} موجب سیانوز می‌شود:
\E{PS}، آترزی \E{P}، آنومالی \E{T}، \E{RV} دیسپلاستیک.
\end{itemize}

\Page{۲}
\begin{itemize}
\item به دلیل موقعیت آناتومیک \E{PFO}، خارج از دورهٔ نوزادی، انتظار شانت \M{L\to R} نداریم.
\item اگر شانت \M{L\to R} وجود داشت، احتمال افزایش فشار \E{LA} در زمینهٔ \E{MS}.
\item یک \E{PFO}، \E{Valve Comp} ولی \E{probe patent} در ۱۵ تا ۳۰\% بزرگسالان دیده می‌شود.
\item \E{PFO} ایزوله نیازمند مداخلهٔ جراحی ندارد، ولی ریسک آمبولی پارادوکسیکال دارد.
\item در بزرگسال با \E{CVA} ترومبوآمبولیک، اکو قلب برای \E{PFO}؛ اگر تأیید شد، بستن \E{PFO} اندیکاسیون دارد.
\end{itemize}

\subsection*{۲. \E{Ostium Secundum Defect (OSD)}}
\begin{itemize}
\item یک \E{OSD} در محل فوسا اووالیس، شایع‌ترین فرم \E{ASD} است.
\item در آن دریچه‌های \E{AV} از نظر ساختار \E{N/L} هستند.
\item با \E{MVP} همراهی دارد ولی اهمیت بالینی ندارد.
\item \E{OSD} ممکن است منفرد یا متعدد (فنستریتهٔ آتریال سپتوم) باشد.
\item \M{2\,\mathrm{cm}<\mathrm{Opening}} در \E{OSD} معمولاً علامت‌دار می‌شود؛ در کودکان به ندرت.
\end{itemize}

\Page{۳}
\begin{itemize}
\item ضایعات بزرگ‌تر ممکن است به سمت \E{inf} گسترش یافته، \E{IVC} را در بر \Cut
\item شاخه‌های شکل دست‌نویس: \E{inf} به \E{IVC}، آتریوم سینوس کرونری؛ \E{Sup} به \E{SVC}؛ \E{Post}.
\item شیوع \E{OSD} در \M{F/M=3/1} است.
\item \E{PAPVR} ممکن است در همراهی با \E{OSD} دیده شود.
\item پارشیال آنومالوس پولمونری ونوس ریترن؛ معمولاً مربوط به \E{right upper P. vein}.
\end{itemize}

شدت شانت \M{L\to R} وابسته به سایز \E{OSD} است و در \E{OSD}های بزرگ‌تر، \M{Q_p:Q_s} معمولاً بین \M{3:1} تا \M{4:1} است.
\begin{latin}
\(Q_p:Q_s\): Pulmonary to systemic blood flow
\end{latin}

در اوایل شیرخوارگی چون بطن راست ضخیم است، شانت \M{L\to R} محدود است. در ادامهٔ شیرخوارگی مقاومت عروق ریوی \E{(PVR)} کاهش می‌یابد \M{\downarrow} و ضخامت \E{RV} نیز کاهش می‌یابد \M{\downarrow}؛ به همین دلیل شدت شانت \M{L\to R} افزایش می‌یابد \M{\uparrow}.

در سیر بیماری، بزرگ شدن \E{RAE} و \E{LVE}، و دیلاتاسیون شریان ریوی.
% The handwritten abbreviation on this line is LVE, retained as written.

\Page{۴}
\begin{itemize}
\item \E{LAE} ممکن است رخ دهد؛ به دلیل افزایش بازگشت وریدی از ریه.
\item \E{LV} و آئورت از نظر اندازه \E{N/L} هستند.
\item علی‌رغم افزایش جریان ریوی، در ابتدا انتظار افزایش فشار شریان ریوی نداریم.
\end{itemize}
\begin{latin}
because of the absence of a high-pressure communication between the pulmonary and systemic circulation.
\end{latin}

\E{PVR} طی کودکی پایین می‌ماند ولی ممکن است در بزرگسالی به تدریج افزایش پیدا کند، تا جایی که شانت ریورس شود و آیزن‌منگر رخ دهد و بیمار سیانوز پیدا کند.

\subsection*{یافته‌های بالینی \E{OSD}}
\begin{itemize}
\item معمولاً کودک بدون علامت است.
\item حتی \E{OSD} بسیار بزرگ هم به ندرت نارسایی قلبی ایجاد می‌کند.
\item ممکن است \E{FTT} خفیف داشته باشد.
\item به دنبال افزایش سن، عدم تحمل ورزش.
\item \E{Mild left precordial bulge}.
\item \E{An RV sys lift may be palpable at the LSB}.
\item \E{Pulmonic ejection click} وجود ندارد.
\item پاتوگنومونیک: \E{S\textsubscript{2}: Widely split and fixed}.
\end{itemize}
در حالت طبیعی طی دم، حجم \E{RV} افزایش می‌یابد و \E{S\textsubscript{2}} اسپلیت می‌شود.

\Page{۵}
\begin{itemize}
\item یک سوفل سیستولیک هوادار، \M{3/6}، مدیوم پیچ، بدون هارش \Cut
\U{که با یک تریل همراه است}؛ بهتر در \E{left middle} و \E{upper} استرنال بوردر.
\item به دلیل جریان افزایش‌یافته از \E{RV outflow} به \U{درون} شریان پولمونری.
\item یک سوفل کوتاه مید دیاستولیک به دلیل افزایش جریان از \E{TV} در \E{left lower}؛ با \E{bell} بهتر شنیده می‌شود.
\item نشان‌دهندهٔ \M{Q_p:Q_s} حداقل \M{2/1}.
\end{itemize}

\subsection*{تشخیص}
\begin{itemize}
\item در \E{CXR}: درجات مختلفی از \E{RVE} و \E{RAE}.
\item افزایش ضخامت شریان پولمونری.
\item افزایش \E{Vasculature} ریه.
\item \U{کاردیوگراف} \E{Enlarged} در پروجکشن لترال بهتر دیده می‌شود.
\item در \E{EKG}: \E{RV volume overload}.
\item \E{RAD or NAX}.
\item \E{\(rsR'\) in right precordial}؛ \E{minor RV conduction delay}.
\item \E{RVH} در نبود \E{PHTN} یا \E{PS} نامعمول است.
\item \E{Echo}: \E{RV volume overload}؛ \E{Abnormal ventricular septum motion}؛ افزایش \E{RV end-D dimension}.
\end{itemize}

\Page{۶}
یک سپتوم \E{N/L} طی \E{S} به سمت \E{post} و طی \E{D} به سمت \E{Ant} حرکت می‌کند؛ \E{Sync} با کانترکشن \E{LV}.

اگر بیمار در \E{Ph/E} و \E{CXR} شواهد \E{ASD} کلاسیک داشته باشد و از نظر همودینامیک علامت‌دار باشد و \E{Echo} صرفاً یک \E{Secundum ASD} ایزوله نشان دهد، می‌تواند مستقیم ترمیم انجام دهد و نیاز نیست کاتتریزیشن تشخیصی انجام دهد؛ به جز بیماران مسن.
\footnote{در اصل، پیش از عبارت همودینامیک، \E{EKG} نوشته و خط زده شده است.}

\subsection*{عوارض \E{Secundum ASD}}
\begin{itemize}
\item معمولاً ایزوله است ولی می‌تواند در همراهی با \E{PAPVR} باشد، یا \E{PS} یا \E{VSD} یا استنوز شاخه‌های شریان پولمونری، یا \E{MVP} یا \E{MVI} یا \E{Left SVC} باقی‌مانده.
\item با سندرم \E{Holt--Oram} همراهی دارد؛ \E{AD}.
\item یک فرم خانوادگی آن با \E{AV conduction delay} همراهی دارد.
\item یک فرم خانوادگی بدون همراهی با بلوک.
\end{itemize}

\subsection*{درمان}
اگر علامت دارد: \E{Catheter} یا \E{Surgery}.

\Page{۷}
\begin{itemize}
\item بدون علامت با \M{Q_p:Q_s} حداقل \M{2:1} یا با \E{RVE}: ترمیم.
\item زمان جراحی الکتیو: از ۱ سالگی تا قبل از مدرسه.
\item مورتالیتی در جراحی قلب باز زیر ۱\% است.
\item بهتر است در کودکی ترمیم شود. در بزرگسال ریسک بالا است \Cut
\item ریسک آریتمی در بزرگسالان بیشتر است \M{\uparrow}.
\item پروسیجر \E{choice}: کاتتر.
\end{itemize}

کدام بیمار نیاز به ترمیم ندارد؟
\begin{latin}
Small Secundum ASD and minimal \(L\to R\) shunt without RVE
\end{latin}

مشخص نیست که آیا یک \E{ASD} کوچک در بزرگسالان ریسک \E{Stroke} را بالا می‌برد یا نه، که نیازمند پروفیلاکسی باشد.

\subsection*{پروگنوز}
\begin{itemize}
\item حتی دفکت \E{moderate} در شیرخوار ممکن است رفع شود \Cut
\item معمولاً تا دههٔ سوم علامت‌دار نمی‌شوند.
\end{itemize}

\Page{۸}
\begin{itemize}
\item علائم در بزرگسال: \E{pHTN}، آریتمی دهلیزی، \E{TI}، \E{MI} (نارسایی میترال)، \E{HF}.
\item در بارداری تشدید علائم رخ می‌دهد.
\item اندوکاردیت عفونی نادر است.
\item \E{ABT} پروفیلاکسی برای \E{ASD} ایزوله جایگاهی ندارد.
\end{itemize}

\subsection*{\E{Sinus Venosus Atrial Septal Defect}}
\begin{itemize}
\item در قسمت فوقانی سپتوم بین‌دهلیزی در نزدیکی \E{SVC}.
\item معمولاً یک یا چند ورید پولمونری به \E{SVC} می‌ریزد.
\item خون وریدی سیستمیک وارد دهلیز چپ می‌شود.
\item به ندرت سیانوز بالینی ایجاد می‌کند.
\item علائم، \E{EKG} و همودینامیک مشابه \E{Secundum ASD}.
\item تشخیص با \E{Echo} است؛ اگر نشد، \E{CT} یا \E{cMRI}.
\item صرفاً در بزرگسالان برای اندازه‌گیری \E{PVR} از \E{Cath} تشخیصی استفاده می‌کنیم.
\end{itemize}

\Page{۹}
\begin{itemize}
\item اصلاح آن معمولاً نیازمند جراحی است.
\item به ندرت ممکن است \E{IVC} را درگیر کند.
\item پروگنوز پس از جراحی عالی است.
\end{itemize}

\subsection*{\E{Partial Anomalous Pulmonary Venous Return (PAPVR)}}
\begin{itemize}
\item یک یا چند ورید پولمونری به اشتباه به \E{SVC} یا \E{IVC}، \E{RA} یا سینوس کرونری می‌ریزند.
\item یک شانت \M{L\to R} ایجاد می‌شود.
\item معمولاً وریدهای یک ریه را درگیر می‌کند؛ معمولاً راست.
\item اگر \E{ASD} وجود داشته باشد، اغلب \E{Sinus Venosus} است ولی می‌تواند \E{Secundum} باشد.
\item در \E{Echo} وقتی \E{ASD} دیدیم باید برای \E{PAPVR} هم فکر کنیم.
\item شرح حال، معاینه، علائم؛ همه‌چیز مثل \E{Ostium Secundum ASD}.
\item بعضاً ورید آنومالوس به \E{IVC} می‌ریزد و در \E{CXR} به صورت یک سایهٔ هلالی یا دانسیتهٔ عروقی در بردر راست قلب دیده می‌شود: \E{Scimitar Syndrome}.
\item در این‌ها \E{ASD} وجود ندارد.
\end{itemize}

\Page{۱۰}
ولی هایپوپلازی ریه شایع است.

\subsection*{\E{Total Anomalous Pulmonary Venous Return (TAPVR)}}
یک ضایعهٔ سیانوتیک.

تشخیص با \E{Echo / MRI / CT}.

پروگنوز \E{PAPVR} عالی است. وقتی یک شانت \M{L\to R} جدی ایجاد شد، جراحی؛ \E{ASD} همراه هم باید بسته شود.

\subsection*{\E{Atrioventricular Septal Defects}}
\begin{latin}
(Ostium Primum, Atrioventricular Canal, Endocardial Cushion Defects)
\end{latin}
\begin{itemize}
\item نقص در سپتوم \E{AV}.
\item در قلب \E{N/L} دریچهٔ \E{T} به اپکس از \U{آن} نزدیک‌تر است.
\item یک سپتوم بین \E{LV} و \E{RA}.
\item اگر کلاً نباشد: \E{Ostium Primum}.
\end{itemize}

\Page{۱۱}
\subsection*{\E{Ostium Primum}}
\begin{itemize}
\item معمولاً دریچهٔ \E{T} عملکرد \E{N/L} دارد.
\item معمولاً یک \E{cleft} روی \E{ant. leaflet} دریچهٔ \E{M} وجود دارد.
\item سپتوم بین‌بطنی سالم است.
\end{itemize}

% Dual-base image path: works standalone (from this folder) and inside the
% notebook build (from the repository root).
\IfFileExists{images/pic01.jpg}{\gdef\CAimagebase{images/}}%
  {\gdef\CAimagebase{pediatrics/ped-cardiology/congenital acyanotic diseases/images/}}
\newcommand{\Original}[2]{%
  \clearpage\subsection*{تصویر #1 ـ اصل دست‌نویس}%
  \begin{center}%
  \includegraphics[width=\linewidth,height=0.88\textheight,keepaspectratio]{\CAimagebase pic#2.jpg}%
  \end{center}}
\Original{۱}{01}
\Original{۲}{02}
\Original{۳}{03}
\Original{۴}{04}
\Original{۵}{05}
\Original{۶}{06}
\Original{۷}{07}
\Original{۸}{08}
\Original{۹}{09}
\Original{۱۰}{10}
\Original{۱۱}{11}
\ifCongenitalAcyanoticStandalone
\end{document}
\fi
